% Authoritative LaTeX; edit this file, not the historical Markdown.
\chapter{From Text to Tokens}\label{chapter-2-from-text-to-tokens}
\section{The model does not receive your
sentence}\label{the-model-does-not-receive-your-sentence}

Suppose two people send what appears to be the same request:

\begin{verbatim}
Explain why token boundaries matter.
\end{verbatim}

One client sends a raw completion string. Another sends a chat message
with the role \texttt{user}. One runtime normalizes the text; another
preserves every byte. One uses the vocabulary packaged with the model;
another accidentally loads a tokenizer with the same family name but a
different revision. All four paths produce valid arrays of integers.
Only some of those arrays are the input the model was trained to
interpret.

The model does not receive a sentence, a chat bubble, or a row of
visible characters. It receives \textbf{token IDs}. Each ID refers to
one entry in one finite, configured \textbf{vocabulary}. The conversion
from text to those IDs is \textbf{tokenization}, and it can include
Unicode normalization, whitespace rules, pre-tokenization, a
segmentation algorithm, byte fallback, and explicit insertion of control
identities. On the output side, generated IDs become byte pieces. Those
pieces may cut through a UTF-8 character, so a runtime sometimes has to
wait for several tokens before it can emit valid text.

This boundary is part of model semantics. If it is wrong, the numerical
model can execute every matrix multiplication correctly and still
receive the wrong sequence.

Chapter 1 made request ownership real while treating the prompt as
opaque. This chapter opens that first black box. We will distinguish
text, Unicode scalar values, bytes, and token IDs; derive byte-pair
encoding (BPE) by hand; separate BPE from SentencePiece's Unigram
algorithm; make special-token insertion explicit; serialize a tiny chat
correctly; and stream decoded bytes without producing malformed UTF-8.
We study the tokenizer independently of numerical inference. The current
executable's generation command uses the four-token model vocabulary, while
its tokenization demonstration uses a separate teaching BPE vocabulary.
An ID from one must never be silently passed to the other.

Chapter 3 will open the numerical function already observed in Chapter 1.
We do not derive that function here.

\section{Four sequences can describe one visible
input}\label{four-sequences-can-describe-one-visible-input}

Begin with a small piece of text:

\begin{verbatim}
世🚀
\end{verbatim}

It is tempting to count two ``characters'' and move on. An inference
engine needs more precise units. The
\href{https://github.com/hermonai/inside-the-llm-engine/blob/astra-visual-rewrite/diagrams/tokenizer/text-unicode-bytes-tokens.txt}{text/Unicode/bytes/tokens
diagram} shows four:

\begin{quote}
\begin{enumerate}
\item visible text       Unicode scalars       UTF-8 bytes           token IDs
\item "世🚀"          \(\longrightarrow\)   U+4E16 U+1F680    \(\longrightarrow\)   E4 B8 96 F0 9F 9A 80  \(\longrightarrow\)  [x, y, z]
\item 2 scalars                7 bytes            model-specific
\end{enumerate}
\end{quote}

The visible text is what a renderer shows. A \textbf{Unicode code point}
is an integer position in the Unicode codespace. A \textbf{Unicode
scalar value} is a code point excluding the surrogate range reserved for
UTF-16. UTF-8 encodes each scalar value as one to four bytes. A
tokenizer then maps a configured transform of those bytes or text units
to vocabulary identifiers.

None of these boundaries has to match another.

The letter \texttt{é} can be one scalar value, U+00E9. It can also be
represented by the canonically equivalent sequence U+0065 LATIN SMALL
LETTER E followed by U+0301 COMBINING ACUTE ACCENT. A user may perceive
one grapheme in both cases, but the second representation has two scalar
values and three UTF-8 bytes. An emoji such as
\texttwemoji{1f469-1f3fd-200d-1f4bb} contains multiple scalar values: a
person, a skin-tone modifier, a zero-width joiner, and a laptop. A
renderer may present the sequence as one grapheme. A tokenizer may split
it into many IDs, sometimes in the middle of one scalar's UTF-8
encoding.

Use \emph{character} only when its meaning is genuinely unambiguous. At
system boundaries, name the unit: grapheme, scalar value, UTF-8 byte, or
token ID.

\begin{quote}
\textbf{FIRST PRINCIPLE} A token is an identity in a particular model
vocabulary. It is not necessarily a word, character, Unicode scalar
value, byte, or output string.
\end{quote}

This distinction immediately explains why token-count rules cannot be
derived from word count. Punctuation may have its own ID or join a
neighboring piece. A leading space may be part of a token. A common word
may be one ID while a rare spelling becomes many. Chinese text has no
obligation to follow space-delimited English boundaries. An emoji can
require several byte-fragment tokens.

\section{UTF-8 is a byte encoding, not
tokenization}\label{utf-8-is-a-byte-encoding-not-tokenization}

Unicode assigns scalar values. UTF-8 specifies how to encode those
values as bytes. Its leading byte determines the expected sequence
length:

\begin{verbatim}
0xxxxxxx                            one byte
110xxxxx 10xxxxxx                  two bytes
1110xxxx 10xxxxxx 10xxxxxx         three bytes
11110xxx 10xxxxxx 10xxxxxx 10xxxxxx four bytes
\end{verbatim}

Those bit shapes are not the whole validity rule. Overlong encodings are
forbidden. UTF-16 surrogate code points are not scalar values and cannot
be encoded as valid UTF-8. Values above U+10FFFF are invalid.
Continuation bytes must appear in the permitted positions. The exact
well-formed ranges are specified by the Unicode Standard and RFC 3629.

Rust's \texttt{str} guarantees valid UTF-8. That is useful for
application text, but a tokenizer/runtime boundary still benefits from
byte-oriented interfaces for three reasons.

First, byte-level vocabularies define pieces in terms of bytes. Second,
generated token pieces may not be valid UTF-8 independently. Third, a
low-level engine needs an explicit policy for arbitrary or malformed
bytes rather than letting a string conversion choose silently.

Our Chapter 2 tokenizer therefore accepts \texttt{\&{[}u8{]}} and
decodes token IDs to \texttt{\&{[}u8{]}}. A separate stream framer
interprets accumulated output as UTF-8. This does not claim that every
tokenizer API should expose raw bytes to every user. It keeps
representation and interpretation separate at the teaching boundary.

\subsection{Malformed bytes require
policy}\label{malformed-bytes-require-policy}

When bytes purport to be UTF-8, Unicode conformance treats ill-formed
sequences as an error condition; they cannot be interpreted as
characters. A product may signal an error, substitute U+FFFD REPLACEMENT
CHARACTER according to a documented maximal-subpart policy, or expose a
binary channel instead of a text channel. What it must not do is
accidentally mix policies.

ENGINE-0 uses a strict policy:

\begin{itemize}
\tightlist
\item
  when the buffer is valid or merely incomplete, emit its longest
  complete valid prefix and retain a suffix that later bytes could
  complete;
\item
  on a definitively invalid append, preserve the buffer and fail without
  emitting a new prefix from that append;
\item
  fail a would-be successful terminal if an incomplete suffix remains;
\item
  never perform lossy replacement.
\end{itemize}

This policy is especially useful for teaching because failure is
observable and raw pending bytes remain available for diagnostics.
Hermon's current production path chooses a different terminal behavior,
which we will inspect later.

\section{A tokenizer is a configured
pipeline}\label{a-tokenizer-is-a-configured-pipeline}

``The tokenizer uses BPE'' is not a complete tokenizer specification. A
useful pipeline model is:

\begin{quote}
\begin{enumerate}
\item input
\item \(\longrightarrow\)  normalization
\item \(\longrightarrow\)  pre-tokenization / boundary rules
\item \(\longrightarrow\)  vocabulary model (BPE, Unigram, WordPiece, ...)
\item \(\longrightarrow\)  special-token post-processing
\item \(\longrightarrow\)  token IDs
\end{enumerate}
\end{quote}

Decoding has its own configured inverse or cleanup steps. Hugging Face
Tokenizers exposes these responsibilities as normalizer, pre-tokenizer,
model, post-processor, and decoder components. Only the model component
is mandatory. That does not make the others unimportant; it means the
artifact must say which ones exist.

\subsection{Normalization}\label{normalization}
\index{tokenization!normalization}

Unicode normalization can replace one valid scalar sequence with an
equivalent or compatibility-related sequence. NFC tends to compose; NFD
decomposes. NFKC and NFKD additionally apply compatibility mappings. A
tokenizer can instead choose identity normalization and preserve the
original sequence.

Other configured normalization may lowercase, strip accents, collapse
whitespace, add a dummy leading space, or replace spaces with a marker.
These are semantic transforms before vocabulary lookup. They can change
token IDs and destroy exact byte round trip.

Consider the composed and decomposed forms of \texttt{é}:

\begin{verbatim}
NFC-like input:  C3 A9
NFD-like input:  65 CC 81
\end{verbatim}

An identity tokenizer sees different bytes. An NFC-normalizing tokenizer
may make them equal before segmentation. Neither behavior is universally
correct. The model and tokenizer configuration determine the contract.

\subsection{Pre-tokenization}\label{pre-tokenization}

A pre-tokenizer introduces boundaries or transforms a stream into
regions the vocabulary model will process. It might separate
punctuation, isolate whitespace, use a regular expression, or map bytes
to a reversible visible alphabet. A BPE model usually does not merge
across a boundary it never sees.

This is one reason two byte-level BPE tokenizers can segment the same
text differently even if both cover all bytes. Their regular
expressions, byte mappings, vocabularies, and merge ranks can differ.

\subsection{Vocabulary model}\label{vocabulary-model}

The model component maps a region into vocabulary pieces and IDs. BPE
applies ranked merge rules. A Unigram model searches segmentations using
learned piece scores. WordPiece uses another greedy vocabulary
procedure. A word-level model may require a pre-tokenizer and map
unknown words to UNK.

The algorithm name still does not give an ID meaning. The exact
vocabulary artifact does.

\subsection{Post-processing and
decoding}\label{post-processing-and-decoding}

A post-processor can add BOS, EOS, separator, or classification tokens
after ordinary encoding. A decoder may reverse a byte mapping, turn a
metaspace symbol back into whitespace, remove continuation markers, or
perform cleanup. Exact round trip depends on the whole chain.

\begin{quote}
\textbf{PROVE IT} For every round-trip claim, state the domain and
configuration. ``All byte sequences round-trip through the toy
byte-fallback BPE with no normalization or special removal'' is
testable. ``Tokenization is reversible'' is too broad.
\end{quote}

\section{Vocabulary identity is model
identity}\label{vocabulary-identity-is-model-identity}

A vocabulary is a finite table. One direction associates IDs with token
definitions; the encoder uses enough additional structure to choose a
sequence of those IDs for input. Let \(V_{\mathrm{vocab}}\) be
vocabulary size and \(N\) be the number of token IDs in one encoded
input. Then every ordinary ID \(x_i\) must satisfy

\begin{equation}
x_i\in\{0,1,\ldots,V_{\mathrm{vocab}}-1\},
\qquad 0\le i<N.
\end{equation}

This is the half-open code interval \texttt{0..V\_vocab} under the
vocabulary's indexing convention. But \texttt{id\ =\ 42} is not
portable. Vocabulary A may define 42 as bytes for a common word
fragment. Vocabulary B may assign 42 to punctuation, a byte fallback
piece, or a control identity.

The vocabulary also does not guarantee that every ID is ordinary text.
Some entries are controls. Some are unused. Some represent unknown
input. Some byte pieces do not form valid text alone.

This is why model loading eventually has to validate more than tensor
shapes. The artifact's weights learned relationships among positions
carrying these exact identities. Swapping a tokenizer changes the
meaning of the embedding row selected for each input position. We will
build embeddings in Chapter 3; for now, retain the contract: stable IDs
must reach that future lookup.

ENGINE-0 introduces a small \texttt{TokenId(u32)} newtype instead of
passing bare integers everywhere. The type cannot prove which vocabulary
an ID belongs to, but it prevents accidental arithmetic and makes
interfaces say when they carry identity rather than text.

\section{Build BPE from first
principles}\label{build-bpe-from-first-principles}
\index{byte-pair encoding}

Byte-pair encoding appears in several tokenizer families, but the
essential encoding operation can be understood with five letters. Our
toy vocabulary starts with all 256 byte values, then adds learned merge
results.

The fixed merge table for \texttt{lower} is:

\begin{quote}
\begin{enumerate}
\item rank 0: l   + o    \(\longrightarrow\)  lo
\item rank 1: lo  + w    \(\longrightarrow\)  low
\item rank 2: e   + r    \(\longrightarrow\)  er
\item rank 3: low + er   \(\longrightarrow\)  lower
\end{enumerate}
\end{quote}

The
\href{https://github.com/hermonai/inside-the-llm-engine/blob/astra-visual-rewrite/diagrams/tokenizer/bpe-merge-process.txt}{BPE
merge diagram} applies it:

\EngineFigure{bpe-merges}{The byte positions stay fixed while token boundaries
disappear. Each highlighted span is the result of applying the lowest-ranked
available merge. A token's width here represents its byte coverage, not a
learned numerical embedding.}{fig:bpe-merges}

The encoder starts from base symbols. It inspects adjacent pairs that
have merge rules, chooses the available rule with the lowest rank,
replaces that pair with the result symbol, and repeats until no rule
applies.

Two details matter.

First, these rules are fixed before the prompt arrives. BPE
\emph{training} learns a merge vocabulary from a corpus. BPE
\emph{encoding} applies the learned artifact. Counting pairs in a user's
prompt and inventing a merge would change the tokenizer during
inference.

Second, order is semantic. The result symbol of a later merge may
require an earlier result. Greedily choosing a visually large piece
without respecting ranks can produce a different segmentation. If the
same rank could apply at multiple positions, the toy contract chooses
the leftmost occurrence, then scans again. Production formats commonly
assign unique ranks to merge pairs, but occurrence order still needs
deterministic behavior.

\subsection{Hand cases}\label{hand-cases}

For \texttt{lolo}, rank 0 applies twice:

\begin{quote}
\begin{enumerate}
\item l | o | l | o
\item [l+o] | l | o  \(\longrightarrow\)  lo | l | o
\item lo | [l+o]      \(\longrightarrow\)  lo | lo
\end{enumerate}
\end{quote}

The result IDs are \texttt{{[}256,\ 256{]}} in our fixture.

For \texttt{xyz}, no merge rule applies. The result remains three base
byte IDs: \texttt{{[}120,\ 121,\ 122{]}}.

For empty input, the result is an empty ID vector. That is a valid
tokenizer result. The generation request layer may independently reject
an empty prompt. Keeping those decisions separate prevents a tokenizer
from owning API policy.

\subsection{A clarity-first
implementation}\label{a-clarity-first-implementation}

\texttt{\seqsplit{TinyBpeTokenizer}} stores the fixed rules and decoded
bytes for every merge result. Construction rejects duplicate pairs,
duplicate ranks, reserved result IDs, duplicate result IDs, and
references to symbols not defined by an earlier rule.

Encoding begins with one \texttt{TokenId} per byte. On every iteration
it scans all adjacent pairs, remembers the lowest-rank rule, merges one
occurrence, and repeats. This is deliberately easy to audit and can be
quadratic in input length. The limitation is part of the lesson. A
production tokenizer may use priority queues, linked symbol structures,
caching, vectorized search, or parallel pre-tokenization, but it must
preserve the configured segmentation.

\begin{quote}
\textbf{BUILD IT} Complete
\href{https://github.com/hermonai/inside-the-llm-engine/blob/astra-visual-rewrite/labs/lab-02-tokenize-by-hand.md}{Lab
2 --- Tokenize by Hand} before changing the merge table. Predict the
states for \texttt{lower}, \texttt{lolo}, and \texttt{xyz}, then compare
with the CLI.
\end{quote}

\section{Byte fallback is coverage, not
one-token-per-byte}\label{byte-fallback-is-coverage-not-one-token-per-byte}

Our toy starts with every byte, so any input byte sequence can be
represented. Known pairs can still merge into larger pieces.
\texttt{blue}, for example, follows \texttt{b+l\ -\textgreater{}\ bl},
\texttt{bl+u\ -\textgreater{}\ blu}, and
\texttt{blu+e\ -\textgreater{}\ blue}, becoming one ID in the toy
vocabulary. Unknown byte patterns remain base IDs.

This gives complete byte coverage without an ordinary unknown path. It
does not mean every input produces one token per byte. It also does not
mean the output is valid UTF-8. The byte sequence \texttt{C3\ 28} can
round-trip exactly through the tokenizer while remaining malformed as
UTF-8.

SentencePiece can optionally include 256 byte fallback pieces.
GPT-2-style byte-level BPE uses a reversible byte-to-Unicode alphabet
before applying BPE. These are related coverage strategies with
different representations. Neither phrase is sufficient to identify a
tokenizer.

OpenAI's GPT-2 reference encoder makes the layers visible. It encodes
input as UTF-8 bytes, maps all bytes to a reversible Unicode alphabet,
splits with a regular expression, applies ranked BPE, and maps pieces to
IDs. Decode performs the reverse vocabulary and byte mapping before
UTF-8 decoding. A token can therefore represent several original bytes,
and its byte fragment can be invalid UTF-8 alone.

\section{SentencePiece names a toolkit, not one
algorithm}\label{sentencepiece-names-a-toolkit-not-one-algorithm}

SentencePiece trains directly from raw sentences and packages
vocabulary, trainer settings, normalization, and other behavior in its
model artifact. It supports several model types. The two important ones
here are BPE and Unigram.

SentencePiece BPE applies a merge-based model. Its surrounding behavior
may use a dummy prefix and a metaspace marker such as \texttt{▁} to
represent word boundaries without assuming that spaces should disappear.

The SentencePiece Unigram model is different. It has a vocabulary of
candidate pieces with learned scores. Encoding searches for a
high-probability segmentation, commonly by dynamic programming. Training
starts with a larger candidate set and prunes pieces according to the
objective. There is no fixed ranked merge trace equivalent to our
\texttt{lower} example.

Calling both ``SentencePiece tokenization'' hides the algorithm
decision. A real engine must inspect the model type stored in the
artifact. It cannot dispatch to BPE merely because a file was produced
by SentencePiece.

\subsection{Unknown tokens and byte
fallback}\label{unknown-tokens-and-byte-fallback}

SentencePiece model configuration includes special pieces such as UNK,
BOS, EOS, and optionally PAD. Character coverage can leave rare input
outside the ordinary alphabet. Without byte fallback, such input can map
to UNK. That mapping loses information: decoding UNK cannot reconstruct
the original bytes.

With byte fallback enabled, out-of-vocabulary scalars can decompose into
UTF-8 byte pieces. This avoids ordinary unknown loss at the segmentation
stage. It does not reverse earlier normalization. If the normalizer
collapsed two spaces to one before fallback, the second space is already
gone.

\section{When does decode(encode(x)) equal
x?}\label{when-does-decodeencodex-equal-x}

Exact round trip is conditional. It can hold when:

\begin{itemize}
\tightlist
\item
  the input domain is covered without UNK loss;
\item
  normalization is identity or otherwise preserves the exact bytes in
  scope;
\item
  pre-tokenization uses a reversible representation;
\item
  decoding reverses every representation transform;
\item
  no special tokens are inserted, removed, or rendered differently;
\item
  no cleanup step changes spaces or punctuation.
\end{itemize}

It can fail when any of those conditions fails.

Our toy BPE has identity normalization, a complete byte alphabet,
reversible merges, and no implicit special insertion. Therefore, over
the domain \(\mathcal{B}^{*}\) of finite byte strings:

\begin{equation}
\operatorname{decode}(\operatorname{encode}(\mathbf{x}))=\mathbf{x},
\qquad \mathbf{x}\in\mathcal{B}^{*}.
\end{equation}

The equality is over bytes. If \texttt{x} is malformed UTF-8, the result
is still the same malformed byte vector and cannot enter the text stream
successfully.

The real SentencePiece comparison later in this chapter demonstrates a
different configured contract. Its byte fallback covers the normalized
input, but its normalizer collapses or removes some whitespace. The
decoded surface is not the original source string for that fixture.

\section{Special tokens are control
identities}\label{special-tokens-are-control-identities}

The canonical
\href{https://github.com/hermonai/inside-the-llm-engine/blob/astra-visual-rewrite/diagrams/tokenizer/special-token-trust-boundary.txt}{special-token
trust boundary} separates untrusted surface text from
template-authorized control identities.

Model vocabularies commonly reserve identities with roles such as:

\begin{itemize}
\tightlist
\item
  \textbf{BOS}: beginning of sequence;
\item
  \textbf{EOS}: end of sequence or one configured end marker;
\item
  \textbf{PAD}: padding used to fill a shape or batch convention;
\item
  \textbf{UNK}: unknown input that could not be represented ordinarily;
\item
  role markers: system, user, assistant, tool, or model-specific roles;
\item
  end-of-turn or separator markers.
\end{itemize}

These names do not define universal IDs. BOS might be automatically
added for one model and forbidden for another. EOS may differ from an
end-of-turn marker. PAD must not automatically be treated as EOS. Some
models have several end-of-generation identities. UNK can be mandatory
even when a modern model rarely reaches it.

Special tokens also cross a trust boundary. Imagine a vocabulary with a
diagnostic spelling
\texttt{\textless{}\textbar{}assistant\textbar{}\textgreater{}}. If the
ordinary encoder automatically recognizes that byte string anywhere,
user content can inject the assistant control identity. The model then
sees a role transition that the structured request did not authorize.

ENGINE-0 prevents that structurally.
\texttt{\seqsplit{Tokenizer::encode(\&[u8])}} encodes ordinary data
only. It has no flag that turns marker-looking text into control IDs.
Trusted template output is a sequence of typed segments:

\begin{verbatim}
pub enum TemplateSegment {
    Text(Vec<u8>),
    Special(SpecialToken),
}
\end{verbatim}

\texttt{Text} goes through ordinary encoding. \texttt{Special} goes
through an explicit special-ID lookup. The string
\texttt{\textless{}\textbar{}assistant\textbar{}\textgreater{}} inside
\texttt{Text} remains its literal bytes. This separation makes the
authority visible in types and tests.

The toy assigns ordinary byte IDs \texttt{0..=255}, merge IDs beginning
at 256, and special IDs beginning at 1000. BOS, EOS, PAD, UNK, SYSTEM,
USER, ASSISTANT, and END\_TURN occupy distinct IDs. An ordinary decode
request for one of those control IDs returns a typed error instead of
pretending the diagnostic name is user text.

\begin{quote}
\textbf{FIRST PRINCIPLE} Ordinary text encoding and special-token
insertion are different operations with different authority.
\end{quote}

\section{A chat template is part of model
input}\label{a-chat-template-is-part-of-model-input}

A chat API presents messages as records:

\begin{verbatim}
role=system, content="Be concise."
role=user,   content="Why does this split?"
\end{verbatim}

A causal language model still receives one ID sequence. A \textbf{chat
template} serializes roles, content, separators, end-of-turn markers,
optional tool data, and the beginning of the generation turn according
to the convention used during training.

The
\href{https://github.com/hermonai/inside-the-llm-engine/blob/astra-visual-rewrite/diagrams/tokenizer/chat-template-pipeline.txt}{chat-template
diagram} shows the boundary:

\begin{quote}
\begin{enumerate}
\item [messages] --render \(\longrightarrow\)  [Text | Special segments] --encode/insert \(\longrightarrow\)  [token IDs]
\end{enumerate}
\end{quote}

Our \texttt{\seqsplit{TinyChatTemplate}} is intentionally smaller than
production template languages:

\begin{verbatim}
BOS
SYSTEM encode(system content) END_TURN
USER   encode(user content)   END_TURN
ASSISTANT                       # when generation should begin
\end{verbatim}

The implementation accepts typed roles, not arbitrary strings. It
returns typed segments, not a marker-filled string that must later be
reparsed. That is enough to demonstrate the semantic owner without
embedding Jinja in the teaching engine.

\subsection{The wrong template can be valid
text}\label{the-wrong-template-can-be-valid-text}

A common fallback is:

\begin{verbatim}
system: Be concise.
user: Why does this split?
\end{verbatim}

Those bytes are valid UTF-8. The tokenizer can encode them. The request
can fit the context. The runtime can execute without error. Yet the IDs
contain no SYSTEM, USER, END\_TURN, or ASSISTANT controls from the tiny
model contract. Validity at the string or array layer does not establish
semantic correctness.

ENGINE-0 cannot measure the effect on answer quality because its model
still returns a fixed candidate table.
\href{https://github.com/hermonai/inside-the-llm-engine/blob/astra-visual-rewrite/labs/lab-04-use-the-wrong-chat-template.md}{Lab
4} therefore compares exact bytes and IDs, not generated quality. A
real-model quality claim would require a pinned model, correct baseline,
workload, and evaluation design.

\subsection{Raw completion is not chat
completion}\label{raw-completion-is-not-chat-completion}

A raw completion surface takes ordinary input intended to continue
directly. A chat completion surface first serializes structured
messages. Even when both contain the same visible user sentence, their
model inputs may differ by BOS, role headers, turn terminators,
assistant prefixes, and whitespace.

Do not automatically apply a chat template to raw input. Do not flatten
chat by intuition. The API surface selects a preparation contract.

\subsection{Avoid duplicate special
insertion}\label{avoid-duplicate-special-insertion}

Hugging Face's official chat-template guidance notes an important
failure mode. A template may already include BOS, EOS, or other
controls. If a caller renders the template to text and then tokenizes
with automatic special addition enabled, it can duplicate them.
Tokenizing through the template-aware surface is often safer because the
owner knows which controls are already present.

\texttt{\seqsplit{add\_generation\_prompt}} is also model-specific. For
templates that use an assistant-turn opener, it appends the sequence
that tells the model an assistant reply comes next. Some templates do
not need such a prefix. The flag does not define a universal string.

\section{Bind model, tokenizer, template, and special
semantics}\label{bind-model-tokenizer-template-and-special-semantics}

A directory containing \texttt{model.gguf}, \texttt{tokenizer.json}, and
a template file is not correct merely because every file parses. The
components must belong to the same model contract.

The
\href{https://github.com/hermonai/inside-the-llm-engine/blob/astra-visual-rewrite/diagrams/tokenizer/model-tokenizer-template-contract.txt}{contract
diagram} makes the dependency visible:

\EngineFigure{token-contract}{Template serialization, tokenizer mapping and model vocabulary rows must agree. Dashed edges denote compatibility constraints, not data movement.}{fig:token-contract}

ENGINE-0 introduces a small \texttt{ModelContract} that records:

\begin{itemize}
\tightlist
\item
  model identity;
\item
  tokenizer name and revision;
\item
  chat-template identity;
\item
  the tokenizer-owned special-ID mapping.
\end{itemize}

Before preparing chat input, it compares the concrete tokenizer and
template identities with the contract. This is not a full artifact
loader. It is a conceptual object that makes one failure explicit:
swapping only one semantic component is not a supported configuration.

Real systems need stronger provenance. A model repository revision,
hashes, embedded metadata, tokenizer assets, template revision, and
runtime support matrix may all participate. Part III will inspect model
containers and metadata. The invariant begins here.

\section{Generated tokens become bytes before they become
text}\label{generated-tokens-become-bytes-before-they-become-text}

On output, a model selects a token ID. The tokenizer vocabulary maps
that ID to a byte piece. Only the concatenation of pieces is expected to
represent output text.

The
\href{https://github.com/hermonai/inside-the-llm-engine/blob/astra-visual-rewrite/diagrams/tokenizer/token-to-byte-stream.txt}{token-to-stream
diagram} separates events and state:

\EngineFigure{utf8-buffer}{Byte buffering is a state transition, not a token-to-word lookup. The second piece completes the first scalar.}{fig:utf8-buffer}

Consider \texttt{世}, whose UTF-8 bytes are \texttt{E4\ B8\ 96}. Suppose
one generated ID maps to \texttt{E4\ B8} and the next maps to
\texttt{96}.

The first selection is a real token event. It advances generation
history and token accounting. It cannot produce a valid text event yet.
The buffer owns two pending bytes. The next token completes the scalar,
after which the runtime can emit \texttt{世} and clear the buffer.

The
\href{https://github.com/hermonai/inside-the-llm-engine/blob/astra-visual-rewrite/diagrams/tokenizer/utf8-partial-token-boundary.txt}{partial-boundary
diagram} shows the timeline:

\begin{quote}
\begin{enumerate}
\item token x bytes: E4 B8        [pending E4 B8]        text: (nothing)
\item token y bytes:       96   \(\longrightarrow\)  [pending E4 B8 96]  \(\longrightarrow\)  validate  \(\longrightarrow\)  text: "世"
\end{enumerate}
\end{quote}

\subsection{The buffer has a small semantic
bound}\label{the-buffer-has-a-small-semantic-bound}

After every complete valid prefix is emitted, a valid incomplete UTF-8
suffix can contain at most three bytes: the beginning of a four-byte
scalar. That does not bound a network output queue or the number of
model tokens. It bounds this specific framing state.

\texttt{\seqsplit{Utf8StreamDecoder::push}} appends one piece and asks
Rust's UTF-8 validator for the status:

\begin{itemize}
\tightlist
\item
  complete valid buffer: emit it and clear;
\item
  error with no definite error length: emit the valid prefix, retain the
  incomplete suffix;
\item
  error with a definite error length: return \texttt{InvalidSequence}
  with the bytes.
\end{itemize}

\texttt{finish} succeeds only when no suffix remains. It never calls
\texttt{\seqsplit{String::from\_utf8\_lossy}}.

\subsection{Token latency and text latency can
differ}\label{token-latency-and-text-latency-can-differ}

Chapter 1 defined time to first token at an observable output boundary.
Now we can name two lower-level events:

\begin{verbatim}
first selected token ready
first valid text piece emitted
\end{verbatim}

They often occur close together for ASCII-heavy output. They can differ
when early token pieces are incomplete UTF-8 or non-rendered controls.
ENGINE-0 records both \texttt{\seqsplit{time\_to\_first\_token}} and
\texttt{\seqsplit{time\_to\_first\_text}}. These trace values remain
instructional, not benchmarks.

\begin{quote}
\textbf{BUILD IT} Complete
\href{https://github.com/hermonai/inside-the-llm-engine/blob/astra-visual-rewrite/labs/lab-03-stream-utf8-across-tokens.md}{Lab
3 --- Stream UTF-8 Across Token Boundaries}. Inject a partial scalar, a
definite malformed sequence, and an incomplete terminal suffix.
\end{quote}

\section{Evolve ENGINE-0 without building a second
runtime}\label{evolve-engine-0-without-building-a-second-runtime}

The Chapter 1 lifecycle remains:

\begin{quote}
\begin{enumerate}
\item validate  \(\longrightarrow\)  admit  \(\longrightarrow\)  execute  \(\longrightarrow\)  select  \(\longrightarrow\)  stream  \(\longrightarrow\)  one terminal outcome
\end{enumerate}
\end{quote}

The input and output boundaries become real:

\begin{quote}
\begin{enumerate}
\item input bytes  \(\longrightarrow\)  tokenizer.encode  \(\longrightarrow\)  Request.input\_tokens
\item numerical model, treated as an opaque boundary in this chapter
\item generated ID  \(\longrightarrow\)  tokenizer.decode\_token  \(\longrightarrow\)  UTF-8 framer  \(\longrightarrow\)  text stream
\end{enumerate}
\end{quote}

The complete ownership map is stored in
\href{https://github.com/hermonai/inside-the-llm-engine/blob/astra-visual-rewrite/diagrams/tokenizer/engine0-token-ownership.txt}{\texttt{\seqsplit{engine0-token-ownership.txt}}}.

\texttt{Request} owns tokenized input rather than an opaque prompt \texttt{String}. It
owns:

\begin{verbatim}
pub struct Request {
    pub id: RequestId,
    pub input_tokens: Vec<TokenId>,
    pub input_bytes: usize,
    pub max_new_tokens: usize,
    pub sampling: SamplingConfig,
}
\end{verbatim}

\texttt{input\_bytes} exists for trace/accounting context. The model
consumes IDs. Whitespace-only input is valid data now; an empty encoded
sequence is rejected by request validation. The tokenizer itself can
still encode empty bytes to an empty vector because that is the
mathematically correct transform.

\texttt{Token} now contains identity and kind, not a cached
\texttt{String}. A historical prototype let its scripted
model own detokenized text; the current tokenizer owns ID-to-byte
meaning.

\texttt{Runtime\textless{}M,\ S,\ T\textgreater{}} owns a model,
selector, and output tokenizer. After an ordinary token is selected, it:

\begin{enumerate}
\def\labelenumi{\arabic{enumi}.}
\tightlist
\item
  appends the identity to runtime-owned generation state;
\item
  emits the token identity event;
\item
  looks up the byte piece;
\item
  pushes the bytes into the request-local UTF-8 framer;
\item
  emits a text event only when a valid non-empty prefix exists;
\item
  continues or reaches one terminal outcome.
\end{enumerate}

EOS remains a control token. It is selected and traced but not decoded
as ordinary output. Before successful EOS or maximum-token completion,
the runtime verifies that the UTF-8 buffer is empty. A pending
incomplete suffix turns the attempted success into
\texttt{\seqsplit{Failed(Utf8Stream(...))}}. Cancellation and an earlier
model failure retain their own terminal causes; pending bytes are not
emitted after a non-success terminal.

\subsection{Keep the two demonstration vocabularies separate}
\label{separate-demonstration-vocabularies}
The BPE probes below use \texttt{TinyBpeTokenizer::teaching}. Generation uses
\texttt{TinyLmTokenizer}, whose four-token vocabulary matches the numerical
fixture. The ID 259 from the BPE probe is not an input to that four-token
model. A historical scripted candidate source is preserved in Git history,
not in the current executable interface described here.

This separation lets us test byte-level tokenization without claiming that a
four-token numerical model understands arbitrary text. The next chapter
derives that numerical fixture, including its deliberate vocabulary limit.

\subsection{CLI probes}\label{cli-probes}

The executable keeps generation and adds two narrow probes:

\begin{verbatim}
cd code/mini-engine
cargo run -p engine0 -- tokenize lower
cargo run -p engine0 -- decode 259
cargo run -p engine0 -- --trace 'I like'
\end{verbatim}

\texttt{tokenize} prints tokenizer identity, byte/scalar/token counts,
IDs, and piece bytes. \texttt{decode} accepts numeric ordinary IDs,
prints each byte piece, and runs the strict UTF-8 framer. The generation
trace adds \texttt{InputEncoded}, \texttt{TokenDecoded},
\texttt{Utf8Buffered}, and \texttt{TextEmitted} stages to the Chapter 1
lifecycle events.

The CLI is a teaching surface, not a standard tokenizer file format or
production protocol.

\section{Prove tokenization
independently}\label{prove-tokenization-independently}

The implementation is not its own oracle. Chapter 2 has three small
committed fixture groups under
\texttt{\seqsplit{code/mini-engine/fixtures/tokenizer/}} and a prose
oracle under \texttt{code/reference/}:

\begin{itemize}
\tightlist
\item
  the BPE rank table and hand encodings;
\item
  byte fragments for valid, malformed, and incomplete UTF-8;
\item
  the tiny chat-template control order.
\end{itemize}

The byte tokenizer is also an independent executable oracle for raw byte
round trips. It maps each byte directly to the numerically equal ID and
does not share the BPE merge algorithm.

The test suite covers four layers.

\subsection{Tokenizer semantics}\label{tokenizer-semantics}

\begin{itemize}
\tightlist
\item
  empty, ASCII, Chinese, emoji, combining, whitespace, repeated, and
  arbitrary malformed byte inputs;
\item
  deterministic encoding and exact byte round trip for the toy BPE;
\item
  merge prerequisites, rank order, leftmost repeat behavior, and
  unmergeable inputs;
\item
  invalid IDs and invalid merge tables;
\item
  ordinary marker-looking text versus explicit special insertion.
\end{itemize}

\subsection{Template and binding
semantics}\label{template-and-binding-semantics}

\begin{itemize}
\tightlist
\item
  exact BOS/role/end-turn/assistant control placement;
\item
  optional generation prompt;
\item
  marker-looking user content remains ordinary;
\item
  naive flattening produces different IDs;
\item
  mismatched tokenizer identity fails before preparation.
\end{itemize}

\subsection{UTF-8 framing semantics}\label{utf-8-framing-semantics}

\begin{itemize}
\tightlist
\item
  partial scalars across token boundaries;
\item
  several complete scalars in one piece;
\item
  valid prefix plus incomplete suffix;
\item
  definite malformed input;
\item
  incomplete terminal suffix;
\item
  empty pieces never create meaningless empty text events.
\end{itemize}

\subsection{Lifecycle integration}\label{lifecycle-integration}

\begin{itemize}
\tightlist
\item
  the \texttt{blue} then EOS oracle remains stable;
\item
  token and text events are separate and ordered;
\item
  cancellation, model failure, tokenizer failure, malformed output, and
  incomplete output each have one terminal;
\item
  no token, text, or trace event follows terminal;
\item
  repeated semantic streams agree with timing excluded.
\end{itemize}

Run the gate:

\begin{verbatim}
cd code/mini-engine
cargo fmt --all -- --check
cargo check --workspace --all-targets
cargo test --workspace
cargo clippy --workspace --all-targets -- -D warnings
\end{verbatim}

\section{Compare two real tokenizers without ranking
them}\label{compare-two-real-tokenizers-without-ranking-them}

Toy fixtures establish correctness, but real tokenizers show how much
the configured artifact matters. The durable experiment record is
\href{https://github.com/hermonai/inside-the-llm-engine/blob/astra-visual-rewrite/research/part-01/tokenizer-comparison.md}{\texttt{\seqsplit{research/part-01/tokenizer-comparison.md}}}.

It compares:

\begin{enumerate}
\def\labelenumi{\arabic{enumi}.}
\tightlist
\item
  OpenAI \texttt{tiktoken} 0.14.0, named encoding \texttt{cl100k\_base};
\item
  SentencePiece 0.2.2 with Google's pinned small Japanese BPE
  byte-fallback test model.
\end{enumerate}

The SentencePiece artifact is intentionally labeled as a test fixture,
not a production LLM tokenizer. The comparison observes counts and
pieces. It makes no speed or quality claim.

{\def\LTcaptype{none} % do not increment counter
\begin{longtable}[]{@{}
  >{\raggedright\arraybackslash}p{(\linewidth - 8\tabcolsep) * \real{0.30}}
  >{\raggedleft\arraybackslash}p{(\linewidth - 8\tabcolsep) * \real{0.14}}
  >{\raggedleft\arraybackslash}p{(\linewidth - 8\tabcolsep) * \real{0.12}}
  >{\raggedleft\arraybackslash}p{(\linewidth - 8\tabcolsep) * \real{0.20}}
  >{\raggedleft\arraybackslash}p{(\linewidth - 8\tabcolsep) * \real{0.24}}@{}}
\toprule\noalign{}
\begin{minipage}[b]{\linewidth}\raggedright
Input
\end{minipage} & \begin{minipage}[b]{\linewidth}\raggedleft
UTF-8 bytes
\end{minipage} & \begin{minipage}[b]{\linewidth}\raggedleft
Scalars
\end{minipage} & \begin{minipage}[b]{\linewidth}\raggedleft
\texttt{cl100k\_base}
\end{minipage} & \begin{minipage}[b]{\linewidth}\raggedleft
SentencePiece fixture
\end{minipage} \\
\midrule\noalign{}
\endhead
\bottomrule\noalign{}
\endlastfoot
\texttt{Tokenizers\ count\ spaces.} & 24 & 24 & 5 & 25 \\
two-line C-like loop & 39 & 39 & 21 & 38 \\
\texttt{模型把文本变成编号。} & 30 & 10 & 10 & 19 \\
\texttwemoji{1f469-1f3fd-200d-1f4bb}\texttwemoji{1f680} & 19 & 5 & 13 &
20 \\
whitespace/newline fixture & 34 & 34 & 9 & 30 \\
\end{longtable}
}

The numbers are not a scoreboard. The small SentencePiece vocabulary was
trained on a different corpus and has different normalization. Its
behavior is useful precisely because the setup is unlike
\texttt{cl100k\_base}.

\subsection{Pieces can cross scalar
boundaries}\label{pieces-can-cross-scalar-boundaries}

For the Chinese input, \texttt{cl100k\_base} splits the UTF-8 bytes of
\texttt{把} across two tokens: one piece contains \texttt{E6\ 8A}; the
next contains \texttt{8A}. Neither is valid UTF-8 alone. The
SentencePiece model uses byte fallback for several Chinese scalars. Both
reconstruct the input after complete byte concatenation.

For \texttwemoji{1f469-1f3fd-200d-1f4bb}\texttwemoji{1f680},
\texttt{cl100k\_base} produces 13 tokens. Many pieces are fragments such
as \texttt{F0\ 9F}, \texttt{91}, and \texttt{A9}. The SentencePiece
fixture produces a dummy metaspace boundary plus one fallback token per
emoji byte, for 20 tokens. This is the production-shaped reason for the
UTF-8 buffer, not an artificial corner case.

\subsection{Normalization can defeat exact surface round
trip}\label{normalization-can-defeat-exact-surface-round-trip}

The whitespace fixture contains leading spaces, a tab, two internal
spaces, a newline, and trailing spaces. \texttt{cl100k\_base} preserves
the exact surface in this experiment. The SentencePiece fixture decodes:

\begin{verbatim}
leading and internal trailing
\end{verbatim}

Its configured normalization removed and collapsed whitespace. Byte
fallback did not restore the original because fallback operated after
normalization. The two-line code input similarly decodes as one
normalized line under this fixture.

This is why a tokenizer contract includes normalization and decoding,
not only the subword algorithm.

\section{Token count is a systems
variable}\label{token-count-is-a-systems-variable}

Tokenization is model semantics, but token count also enters system
economics. If an encoded prompt has length \(N\), later stages see \(N\)
positions, not the number of words or bytes the caller expected.

For a model context capacity \(T_{\max}\) and requested generation
budget \(N_{\mathrm{new}}\), admission must at least respect the
dimensional constraint

\begin{equation}
N+N_{\mathrm{new}}\le T_{\max}\quad[\mathrm{token\ positions}].
\end{equation}

Token count affects:

\begin{itemize}
\tightlist
\item
  whether input fits the model's context limit;
\item
  how much prefill work Chapter 20 will analyze;
\item
  how much per-position KV state later chapters allocate;
\item
  prefix/cache matching in token space;
\item
  maximum-generation and stop accounting;
\item
  batching shapes and admission budgets;
\item
  provider billing when tokens are the accounting unit;
\item
  rate limits and observability labels.
\end{itemize}

Do not turn this list into a premature performance formula. Different
model architectures, batching, hardware, and token distributions change
costs. The durable statement is dimensional: the engine schedules token
positions, so the correct tokenizer's \(N\) is an input to later work
and memory models. The canonical
\href{https://github.com/hermonai/inside-the-llm-engine/blob/astra-visual-rewrite/diagrams/tokenizer/token-count-system-impact.txt}{system-impact
map} shows those downstream dependencies without inventing a cost
coefficient.

A tokenizer with fewer tokens for one sentence is not automatically
better. Vocabulary size \texttt{V} also affects future embedding and
output dimensions. A larger vocabulary can reduce sequence length for
some inputs while increasing parameter rows and changing training
behavior. Language coverage and segmentation quality matter. We have not
measured any end-to-end tradeoff here.

\subsection{Cache boundaries use IDs, not reconstructed
text}\label{cache-boundaries-use-ids-not-reconstructed-text}

Two text surfaces that render identically after normalization may have
different original bytes. Two tokenizers can produce different ID
sequences for the same bytes. Prefix reuse in a model runtime must
compare the identities used for model state, not a casually re-rendered
string.

Hermon's current batched runtime tokenizes before choosing its best
prefix slot because the reusable state corresponds to token positions.
Its separate \texttt{\seqsplit{hermon-tokenizer}} crate sketches a
tokenizer prefix cache, but that crate does not own the current
real-model route.

\section{Inside Hermon: pinned llama.cpp owns the current
path}\label{inside-hermon-pinned-llama.cpp-owns-the-current-path}

The following account is \textbf{CURRENT} for Hermon commit
\href{https://github.com/hermonai/hermon/commit/472a44cdb511b2dae6c9569e59543db8f8350b25}{\texttt{472a44c}},
inspected on 2026-09-02. The detailed evidence is in the
\href{https://github.com/hermonai/inside-the-llm-engine/blob/astra-visual-rewrite/research/part-01/chapter-02-from-text-to-tokens.md}{Chapter
2 research note}.

The Hermon repository pins llama.cpp submodule commit
\href{https://github.com/ggml-org/llama.cpp/commit/389ff61d77b5c71cec0cf92fe4e5d01ace80b797}{\texttt{389ff61}}.
The upstream llama.cpp head observed during the inspection was newer.
Pinned does not mean stale by itself, and CURRENT does not mean ``equal
to upstream HEAD.'' It means that this pinned source owns the reachable
production path we traced.

\begin{quote}
\textbf{INSIDE HERMON --- CURRENT} The default \texttt{BatchedRuntime}
constructs a message view, asks the llama.cpp model wrapper to apply the
model's embedded chat template with an assistant generation prompt, and
uses a naive role flattener only when no embedded template exists. It
then calls context tokenization with special addition and special
parsing enabled.
\end{quote}

The Rust safe wrapper in
\texttt{\seqsplit{hermon-llamacpp/src/linked.rs}} crosses a narrow C
shim. The shim gets the model's llama.cpp vocabulary and calls
\texttt{llama\_tokenize}. The two Boolean controls matter: llama.cpp
documents \texttt{add\_special} as allowing configured BOS/EOS addition
and \texttt{parse\_special} as allowing control-token spellings to be
recognized instead of treated as plaintext.

Hermon's blocking raw stream calls the same tokenizer with
\texttt{\seqsplit{parse\_special=false}}. Its chat path first applies a
trusted template, then uses true. This is the same authority distinction
our typed teaching segments make, expressed through the pinned upstream
API.

\begin{quote}
\textbf{INSIDE HERMON --- LIBRARY}
\texttt{\seqsplit{crates/hermon-tokenizer}} currently defines a
tokenizer-kind enum, a \texttt{Tokenize} trait, and a prefix-cache
skeleton. Its own source labels BPE, SentencePiece, Tiktoken, and
Hugging Face ingestion as future work. Repository search found no
current real-model request route using it.
\end{quote}

Calling that crate ``Hermon's tokenizer'' without the status would be
misleading. The real-model owner today is the pinned llama.cpp
vocabulary path. The crate is a LIBRARY/research surface.

\subsection{Hermon's UTF-8 buffer}\label{hermons-utf-8-buffer}

Each active batched sequence owns
\texttt{utf8\_buf:\ Vec\textless{}u8\textgreater{}}. After sampling, the
worker calls \texttt{token\_to\_piece(token,\ false)}, appends the
resulting bytes, finds the longest prefix accepted by Rust's UTF-8
validator, emits that prefix as a
\texttt{\seqsplit{StreamItem::Piece(String)}}, and retains the rest.
This is direct production evidence that piece and token boundaries
differ.

The current implementation has a bounded caveat.

\begin{quote}
\textbf{INSIDE HERMON --- CURRENT CAVEAT} The buffer path uses
\texttt{valid\_up\_to()} without separately testing whether
\texttt{error\_len()} identifies a definite malformed sequence. At
successful finalization it converts leftover bytes with
\texttt{\seqsplit{String::from\_utf8\_lossy}}.
\end{quote}

Therefore Hermon's current policy does not match ENGINE-0's strict error
policy. A malformed suffix can remain buffered, and an incomplete or
invalid terminal suffix is rendered with replacement rather than
becoming an explicit runtime error. The blocking convenience engine has
a similar lossy terminal flush.

This chapter records the behavior; it does not modify Hermon. Later
production streaming and protocol chapters must decide the desired
policy, test it at the runtime boundary, and trace terminal cause to the
wire.

\subsection{What tests establish}\label{what-tests-establish}

Hermon's ignored real-model tests exercise tokenization, chat-template
application, token-to-piece concatenation, and next-token differential
behavior when an external model fixture and linked backend are supplied.
They establish useful model-dependent evidence when run. Ordinary CI
without the fixture does not prove every tokenizer/model combination.

The inspected default path has no focused unit test separating
valid-incomplete from definitively malformed UTF-8 fragments. ENGINE-0's
tests are evidence for the teaching contract, not retroactive proof of
Hermon.

\section{Common mistakes}\label{common-mistakes-1}

\subsection{``A token is a word''}\label{a-token-is-a-word}

Tokens can contain a whole word, a prefix, punctuation, a leading space,
many bytes, or an incomplete scalar fragment. Use vocabulary identity as
the definition.

\subsection{``UTF-8 characters are
tokens''}\label{utf-8-characters-are-tokens}

UTF-8 encodes scalar values into bytes. Tokenization is a separate
configured mapping. A tokenizer can split within the byte encoding of
one scalar.

\subsection{``SentencePiece means
Unigram''}\label{sentencepiece-means-unigram}

SentencePiece supports Unigram, BPE, word, and character model types.
Inspect the serialized model configuration.

\subsection{``Byte fallback guarantees exact input round
trip''}\label{byte-fallback-guarantees-exact-input-round-trip}

It guarantees coverage of bytes presented to the vocabulary model.
Earlier normalization can already have changed whitespace or scalar
sequences.

\subsection{``Decode each token as a
string''}\label{decode-each-token-as-a-string}

An individual piece may be invalid UTF-8. Decode to bytes, buffer per
request, and emit only complete valid prefixes under a named
malformed-data policy.

\subsection{``The marker string is the special
token''}\label{the-marker-string-is-the-special-token}

The control identity is the token. A printed marker is a diagnostic
surface. Ordinary user bytes that resemble it must not gain control
authority implicitly.

\subsection{``Any chat formatting is close
enough''}\label{any-chat-formatting-is-close-enough}

The template defines exact model input. A plausible
\texttt{role:\ content} string can be syntactically valid and
semantically wrong.

\subsection{``Add BOS and EOS everywhere for
safety''}\label{add-bos-and-eos-everywhere-for-safety}

Templates and tokenizer post-processors may already insert required
controls. Duplication changes the sequence. Follow the bound model
contract.

\subsection{``Fewer tokens is always
better''}\label{fewer-tokens-is-always-better}

Token count affects work, but vocabulary size, language coverage,
learned semantics, model compatibility, and workload also matter. Counts
from unrelated tokenizers are not a quality leaderboard.

\begin{quote}
\textbf{ENGINEERING FAILURE} A service loads weights from revision A and
a tokenizer/template from revision B. Every file parses and all token
IDs are in range. The model emits poor or strangely terminated output.
The defect is not numerical: the same integers select different learned
vocabulary rows and the role controls no longer match training. Artifact
identity must bind the components before execution.
\end{quote}

\section{Exercises}\label{exercises-1}

\subsection{CHECK}\label{check-1}

\begin{enumerate}
\def\labelenumi{\arabic{enumi}.}
\tightlist
\item
  For \texttt{é} written as U+0065 U+0301, count visible graphemes,
  scalar values, UTF-8 bytes, and tokens under the toy BPE. Repeat after
  NFC composition.
\item
  Explain why byte vectors can round-trip through
  \texttt{\seqsplit{TinyBpeTokenizer}} even when they cannot enter
  \texttt{\seqsplit{Utf8StreamDecoder}} successfully.
\item
  Classify BOS, EOS, PAD, UNK, USER, and END\_TURN as ordinary content,
  controls, or information-loss markers. Which names can share behavior
  only if model metadata says so?
\item
  A prompt has 1,000 Unicode scalar values and 1,600 tokens. Which count
  enters a 1,200-token context limit? What additional model/runtime
  facts are needed to estimate memory or latency?
\end{enumerate}

\subsection{BUILD}\label{build-1}

Complete
\href{https://github.com/hermonai/inside-the-llm-engine/blob/astra-visual-rewrite/labs/lab-02-tokenize-by-hand.md}{Lab
2},
\href{https://github.com/hermonai/inside-the-llm-engine/blob/astra-visual-rewrite/labs/lab-03-stream-utf8-across-tokens.md}{Lab
3}, and
\href{https://github.com/hermonai/inside-the-llm-engine/blob/astra-visual-rewrite/labs/lab-04-use-the-wrong-chat-template.md}{Lab
4}. Keep the hand oracle, implementation, and failure injection
separate.

Add a CLI input containing a tab, newline, Chinese text, and an emoji.
Record bytes, scalar values, IDs, and decoded bytes. Do not call the ID
count a word count.

\subsection{BREAK}\label{break-1}

In a temporary branch:

\begin{enumerate}
\def\labelenumi{\arabic{enumi}.}
\tightlist
\item
  make ordinary encoding recognize
  \texttt{\textless{}\textbar{}assistant\textbar{}\textgreater{}} as
  control ID 1006;
\item
  feed the marker inside a user message;
\item
  observe the special-insertion test fail;
\item
  restore the typed boundary.
\end{enumerate}

Then replace the strict terminal check with lossy UTF-8 conversion. Run
the malformed and incomplete fixtures and explain which evidence
disappeared.

\subsection{EXTEND}\label{extend-1}

Implement an alternative tiny tokenizer configuration that applies one
named normalization rule before the same BPE. Add paired
composed/decomposed and whitespace fixtures. State exactly which byte
round trips no longer hold.

Do not add a third-party dependency merely to hide the transform. The
purpose is to make the semantic difference inspectable.

\section{What this chapter has not
built}\label{what-this-chapter-has-not-built}

This chapter has not derived the numerical language model. Tokenization
establishes categorical identities and byte semantics, not the embedding
values or the arithmetic that produces logits. The current executable
contains those later components, but their derivation belongs to Chapter 3.

We also have not built a production tokenizer loader, trained a
vocabulary, implemented efficient BPE, benchmarked allocations, or
selected the real model artifact for later parts. The real-tokenizer
experiment uses external packages in a temporary environment and commits
only small textual observations.

We have not derived stochastic sampling, temperature, softmax, top-k, top-p
or random-number state here. Chapter 4 opens those components after genuine
logits are understood. Stop-string matching remains outside the implemented
teaching runtime.

\section{Summary}\label{summary-1}

The visible input to an application passes through several
representations before a model can use it. Graphemes, Unicode scalar
values, UTF-8 bytes, and token IDs are different sequences with
different boundaries. A token is one identity in one configured
vocabulary.

A tokenizer is more than a segmentation algorithm. Normalization,
pre-tokenization, vocabulary/model rules, post-processing, special
tokens, and decoding jointly determine behavior. BPE encoding applies a
fixed ranked merge artifact; it does not learn from the prompt.
SentencePiece can store BPE or Unigram models, and the algorithms are
not interchangeable. Byte fallback provides coverage after configured
preprocessing but cannot recover information already removed by
normalization.

Special tokens are controls with explicit authority. Ordinary text that
looks like a marker must remain ordinary. Chat templates serialize
structured messages into the exact sequence expected by the model,
including role and turn boundaries. Raw completion and chat completion
therefore need different preparation paths. Model, tokenizer revision,
template, and special semantics form one contract.

On output, vocabulary lookup produces bytes. Token pieces can split
UTF-8 scalar values, so a per-request decoder buffers a valid incomplete
suffix and emits only complete text. Token identity events and
text-piece events are not one-to-one. Malformed and incomplete terminal
behavior must be a named policy.

ENGINE-0 now implements these boundaries with no external Rust
dependencies. Its request owns token IDs, its tokenizer owns ID-to-byte
meaning, its template inserts typed controls, its contract rejects
mismatched identities, and its UTF-8 framer rejects malformed output
without lossy replacement. Its request lifecycle and exactly-once
terminal behavior remain intact.

Hermon's current real-model path reaches the pinned llama.cpp tokenizer
and embedded chat-template APIs. The separate
\texttt{\seqsplit{hermon-tokenizer}} crate remains a LIBRARY skeleton.
Hermon's per-sequence UTF-8 buffer demonstrates the real token/piece
mismatch, while its lossy terminal flush documents a policy difference
for later audit.

\section{Next: the smallest possible language
model}\label{next-the-smallest-possible-language-model}

Chapter 3 receives stable token IDs and derives the small numerical language
model already exercised by the current CLI. We will distinguish immutable
model parameters from runtime
activations, look up an embedding row, carry a hidden vector through an
output projection, produce one logit per vocabulary item, account for
every shape, and verify the result with an independent hand-computable
oracle.

\section{Primary references}\label{primary-references-1}

\begin{itemize}
\tightlist
\item
  The Unicode Consortium,
  \href{https://www.unicode.org/versions/Unicode17.0.0/core-spec/chapter-3/}{The
  Unicode Standard, Version 17.0.0, Chapter 3} and
  \href{https://www.unicode.org/reports/tr15/}{Unicode Standard Annex
  \#15: Unicode Normalization Forms}.
\item
  IETF, \href{https://www.rfc-editor.org/rfc/rfc3629}{RFC 3629 ---
  UTF-8, a transformation format of ISO 10646}.
\item
  Rico Sennrich, Barry Haddow, and Alexandra Birch, ``Neural Machine
  Translation of Rare Words with Subword Units,'' and the official
  \href{https://github.com/rsennrich/subword-nmt/tree/92d6139d07d30e12735a0af9e7f7f925ebe62c54}{\texttt{subword-nmt}
  implementation at \texttt{92d6139}}.
\item
  OpenAI,
  \href{https://github.com/openai/gpt-2/blob/9b63575ef42771a015060c964af2c3da4cf7c8ab/src/encoder.py}{GPT-2
  \texttt{encoder.py} at \texttt{9b63575}} and
  \href{https://github.com/openai/tiktoken}{\texttt{tiktoken}}.
\item
  Google,
  \href{https://github.com/google/sentencepiece/tree/ac0f71dcbc85f94292266e55bf6cee1b4d6c9dc1}{SentencePiece
  source and documentation at \texttt{ac0f71d}}, including the
  serialized model type and normalization/special-symbol docs.
\item
  Hugging Face,
  \href{https://huggingface.co/docs/tokenizers/components}{Tokenizer
  components} and
  \href{https://huggingface.co/docs/transformers/chat_templating}{chat-template
  guidance}, source snapshots \texttt{d582781} and \texttt{e15d467}
  recorded in the research note.
\item
  llama.cpp,
  \href{https://github.com/ggml-org/llama.cpp/blob/389ff61d77b5c71cec0cf92fe4e5d01ace80b797/include/llama.h}{tokenization
  and chat-template API at Hermon's pinned \texttt{389ff61}}.
\item
  Hermon,
  \href{https://github.com/hermonai/hermon/blob/472a44cdb511b2dae6c9569e59543db8f8350b25/crates/hermon-runtime/src/batched.rs}{\texttt{batched.rs}},
  \href{https://github.com/hermonai/hermon/blob/472a44cdb511b2dae6c9569e59543db8f8350b25/crates/hermon-llamacpp/src/linked.rs}{\texttt{hermon-llamacpp}},
  and
  \href{https://github.com/hermonai/hermon/blob/472a44cdb511b2dae6c9569e59543db8f8350b25/crates/hermon-tokenizer/src/lib.rs}{\texttt{\seqsplit{hermon-tokenizer}}}
  at commit \texttt{472a44c}.
\end{itemize}


\section{Worked synthesis: identities, bytes and evidence}
\subsection{A vocabulary permutation that passes shape checks}
\textbf{Problem.} Two tokenizers have the same vocabulary size. One exchanges
the IDs for two ordinary tokens. Why can model shape validation succeed while
the answer changes?

\textbf{Solution.} A shape check proves only that an ID is in range. The same
integer now selects a different embedding row and labels a different output
logit. If only the tokenizer changes, the model's learned association between
row and token is broken. A coordinated permutation of embedding rows,
output rows, biases and all special-token semantics could preserve behavior,
but changing one file does not perform that transformation.

\subsection{Complete the byte trace}
\textbf{Problem.} A token emits \texttt{E4 B8}; the next emits \texttt{96}.
List pending bytes and text after each token. What changes if EOS arrives
instead of the second token?

\textbf{Solution.} After the first token, pending bytes are
\texttt{E4 B8} and text is empty. After the second, the pending buffer is empty
and U+4E16 has been emitted once. EOS after the first token leaves an incomplete
UTF-8 scalar; under this engine's strict completion policy it is a decoding
failure, not successful text and not an implicit replacement character.

\subsection{Build a discriminating tokenizer check}
\textbf{Problem.} Is encode/decode round-trip equality sufficient to validate
compatibility with a model?

\textbf{Solution.} No. A self-consistent but permuted vocabulary round-trips its
own IDs. Add known byte-to-ID fixtures from the intended model revision,
special-token checks and template-rendering examples. Round trips test internal
consistency; known fixtures test the external semantic contract.

